\begingroup
\renewcommand{\arraystretch}{1.4}
\SourceLeaf{312}
\InterpolationArticle{25}
Si functio $X$ cum terminis $\cos mx$, $\sin mx$ non abrumpitur, sed ulterius excurrit, coëfficientesque terminorum sequentium adhuc nimis considerabiles sunt, ita ut ipsos negligere non liceat, coëfficientes $\gamma,\gamma',\delta'$ etc. in functione $X'$, quae $\mu$ valoribus functionis $X$ periodum completam formantibus, pro $x=a,b,c,d$ etc. satisfacit, notabiliter a coëfficientibus respondentibus $\alpha,\alpha',\beta'$ etc. discrepabunt. Quodsi itaque pro periodo $x=a,b,c,d$ etc. periodum aliam $\mu$ terminorum, puta pro $x=a',b',c',d'$ etc. simili modo tractamus, functionemque $m^{\mathrm{ti}}$ ordinis, his novis valoribus functionis $X$ satisfacientem evolvimus, et perinde ut $X'$ per expressionem
\[
\begin{aligned}
&\gamma+\gamma'\cos x+\gamma''\cos2x+.\,.\,.+\gamma^m\cos mx\\[5pt]
&\phantom{\gamma}+\delta'\sin x+\delta''\sin2x+.\,.\,.+\delta^m\sin mx
\end{aligned}
\]
exhibemus: pro coëfficientibus $\gamma,\gamma',\delta'$ etc. iam valores nanciscemur, ab iis, quos ante invenimus, notabiliter diversos. Hinc intelligitur, quo pacto hos coëfficientes tamquam \textit{variabiles} spectare liceat, et quidem, quum ex artt. 21, 23 singulos perinde per arcum $\mu a'$ determinari pateat, ut valores priores per arcum $\mu a$, singuli coëfficientes considerari poterunt tamquam functiones arcus indeterminati $y$, similiter ut $X$ est functio arcus $x$. Iam supponamus, $\mu\nu=\pi$ valores functionis $X$ periodum integram formantes, pro
\[
\begin{aligned}
&x=a,a',a'',.\,.\,.b,b',b'',.\,.\,.c,c',c'',.\,.\,.d,d',d''.\,.\,.\text{ etc.}
\end{aligned}
\]
cognitos esse, ita ut sit
\[
\begin{aligned}
&a'=a+\tfrac1\pi360^0.\\[5pt]
&a''=a+\tfrac2\pi360^0\text{ etc.}\\[5pt]
&b=a+\tfrac\nu\pi360^0=a+\tfrac1\mu360^0,\\[5pt]
&b'=a+\tfrac{\nu+1}\pi360^0=a'+\tfrac1\mu360^0\text{ etc.}\\[5pt]
&c=a+\tfrac{2\nu}\pi360^0=a+\tfrac2\mu360^0\text{ etc.}
\end{aligned}
\]
pateatque, hanc periodum $\pi$ terminorum in $\nu$ periodos $\mu$ terminorum pro
\[
\begin{aligned}
&x=a,\ b,\ c,\ d\text{ etc.}\\[5pt]
&x=a',\ b',\ c',\ d'\text{ etc.}\\[5pt]
&x=a'',\ b'',\ c'',\ d''\text{ etc.}\\[5pt]
&\text{etc.}
\end{aligned}
\]
discerpi posse. Quodsi itaque eo quo docuimus modo pro singulis periodis functio
\SourceLeaf{313}
$X'$ evolvitur, nanciscimur $\nu$ valores singulorum coëfficientium $\gamma,\gamma',\delta'$ etc., qui hisce deinceps valoribus ipsius $y$ respondent
\[
\begin{aligned}
&y=\mu a,\\[5pt]
& y=\mu a'=\mu a+\tfrac1\nu360^0,\\[5pt]
& y=\mu a''=\mu a+\tfrac2\nu360^0\text{ etc.}
\end{aligned}
\]
\textit{adeoque periodos completas $\nu$ terminorum formant}. Quocirca, per methodum nostram, singuli coëfficientes $\gamma,\gamma',\delta'$ etc. per expressionem talem
\[
\begin{aligned}
&\varepsilon+\varepsilon'\cos y+\varepsilon''\cos2y+.\,.\,.\,.+\varepsilon^n\cos ny\\[5pt]
&\phantom{\varepsilon}+\zeta'\sin y+\zeta''\sin2y+.\,.\,.\,.+\zeta^n\sin ny
\end{aligned}
\]
exhiberi poterunt, ita ut sit $n$ vel $=\tfrac12\nu-\tfrac12$ vel $=\tfrac12\nu$, prout $\nu$ impar est vel par. Manifestum est, talem expressionem eundem valorem ex substitutione $y=\mu b,\mu c,\mu d$ etc. obtinere, quem obtinet ex $y=\mu a$; nec non eundem ex $y=\mu b',\mu c',\mu d'$ etc., quem obtinet ex $y=\mu a'$ et sic porro. Quamobrem si in
\[
\begin{aligned}
&\gamma+\gamma'\cos x+\gamma''\cos2x+\text{ etc.}\\[5pt]
&\phantom{\gamma}+\delta'\sin x+\delta''\sin2x+\text{ etc.}
\end{aligned}
\]
pro singulis coëfficientibus $\gamma,\gamma',\delta'$ etc. hae functiones arcus indeterminati $y$ substituuntur, prodit functio $Z$, duos arcus $x,y$ involvens, transibitque $Z$ in eam functionem $X'$, quae e periodo prima deducta est, si pro $y$ aliquis valorum $\mu a,\mu b,\mu c,\mu d$ etc. substituitur; si vero pro $y$ aliquis valorum $\mu a',\mu b',\mu c',\mu d'$ etc. accipitur, transibit $Z$ in eam functionem, quae eruta est e periodo secunda. Hinc tandem colligitur, si in $Z$ pro $y$ statim scribatur $\mu x$, ita ut functio solius $x$ prodeat ($X''$), hanc ita comparatam fore, ut pro omnibus $\pi$ valoribus ipsius $x$, valoribus propositis functionis $X$ satisfaciat.

\InterpolationArticle{26}
Si methodus artt. 20, 22 ad omnes $\pi$ valores propositos immediate sine praevia discerptione in $\nu$ periodos minores applicata esset, prodiisset functio ($X'''$) ordinis $\tfrac12\pi-\tfrac12^{\mathrm{ti}}$ vel $\tfrac12\pi^{\mathrm{ti}}$, prout $\pi$ impar vel par, quam cum functione $X''$ iam comparabimus. Quum haec posterior, si rite reducta est, manifesto ad ordinem $\mu n+m^{\mathrm{tum}}$ ascendat, tres casus distinguemus.

I. Quando tum $\mu$ tum $\nu$ est impar, erit $\mu n+m=\tfrac12\mu\nu-\tfrac12=\tfrac12\pi-\tfrac12$. unde patet, $X''$ cum $X'''$ prorsus identicam esse debere.

II. Quando $\mu$ par est, $\nu$ impar, fit $\mu n+m=\tfrac12\mu\nu=\tfrac12\pi$, adeoque $X''$ eius\SourceLeaf{314}%
dem ordinis ut $X'''$. Sufficit itaque utriusque functionis terminos ultimos comparare. Haud difficile perspicitur, hos terminos ultimos in $X''$ prodire unice ex terminis $\gamma^m\cos mx+\delta^m\sin mx$, et quidem ex iis terminis coëfficientium $\gamma^m,\delta^m$, qui continent $\cos ny$ et $\sin ny$ sive $\cos\mu nx$ et $\sin\mu nx$. Sint hi termini
\[
\begin{aligned}
&\text{in }\gamma^m\quad.\,.\,.\,.\quad k\cos\mu nx+l\sin\mu nx\\[5pt]
&\text{in }\delta^m\quad.\,.\,.\,.\quad k'\cos\mu nx+l'\sin\mu nx
\end{aligned}
\]
produceturque hinc in $X''$
\[
\begin{aligned}
&(k\cos\mu nx+l\sin\mu nx)\cos mx\\[5pt]
&\qquad{}+(k'\cos\mu nx+l'\sin\mu nx)\sin mx
\end{aligned}
\]
unde prodeunt termini ordinis $\mu n+m^{\mathrm{ti}}$ sequentes
\[
\begin{aligned}
&\tfrac12(k-l')\cos(\mu n+m)x+\tfrac12(l+k')\sin(\mu n+m)x
\end{aligned}
\]
Iam haud difficile quidem \textit{ex formatione coëfficientium} demonstrari potest, esse
\[
\begin{aligned}
&(k-l')\sin(\mu n+m)a=(l+k')\cos(\mu n+m)a
\end{aligned}
\]
praeferimus tamen methodum sequentem, quae concinnior videtur. Ex art. 24 sequitur, esse
\[
\begin{aligned}
&\gamma^m\sin\tfrac12y-\delta^m\cos\tfrac12y=0
\end{aligned}
\]
\textit{tum} si pro $y$ substituitur $\mu a$, atque pro $\gamma^m,\delta^m$ valores ex periodo prima deducti;
\textit{tum} si pro $y$ substituitur $\mu a'$, atque pro $\gamma^m,\delta^m$ valores ex periodo secunda deducti etc. Quodsi itaque pro $\gamma^m,\delta^m$ expressiones, quae erutae sunt indefinitae substituuntur, ita ut $\gamma^m\sin\tfrac12y-\delta^m\cos\tfrac12y$ fiat functio arcus $y$ formae $G$ ad ordinem $n+\tfrac12=\tfrac12\nu$ ascendens, haec pro omnibus $\nu$ valoribus ipsius $y$ his $\mu a,\mu a',\mu a''$ etc. fiet $=0$, adeoque per art. 14 per
\[
\begin{aligned}
&\sin\tfrac12(y-\mu a)\sin\tfrac12(y-\mu a')\sin\tfrac12(y-\mu a'')\text{ etc.}
\end{aligned}
\]
divisibilis, hinc etiam per $\sin\tfrac12\nu(y-\mu a)$ et proin formae $h\sin\tfrac12\nu(y-\mu a)$, ita ut $h$ sit quantitas definita. Quamobrem termini ultimi functionis $\gamma^m\sin\tfrac12y-\delta^m\cos\tfrac12y$ evolutae esse debebunt
\[
\begin{aligned}
&-h\sin\tfrac12\pi a\cos\tfrac12\nu y+h\cos\tfrac12\pi a\sin\tfrac12\nu y
\end{aligned}
\]
omnes vero antecedentes evanescere: illos vero fieri patet
\SourceLeaf{315}
\[
\begin{aligned}
&=-\tfrac12(l+k')\cos(n+\tfrac12)y+\tfrac12(k-l')\sin(n+\tfrac12)y
\end{aligned}
\]
Quocirca erit
\[
\begin{aligned}
&\tfrac12(k-l')=h\cos\tfrac12\pi a=h\cos\tfrac12(\mu n+m)a\\[5pt]
&\tfrac12(l+k')=h\sin\tfrac12\pi a=h\sin\tfrac12(\mu n+m)a\qquad\text{Q.~E.~D.}
\end{aligned}
\]
Haec vero aequatio ipsissimum erit criterium coëfficientium ultimorum functionis
$X'''$, ut in art. 24 ostentum est: quare in hoc quoque casu functio $X''$ cum $X'''$
prorsus identica erit.

III. Quando $\nu$ par est, fit $\mu n+m=\tfrac12\pi+m$, adeoque $X''$ ordinis altioris
quam $X'''$. Facile autem perspicitur, omnes terminos in $X''$, qui ordinem
$\tfrac12\pi^{\mathrm{tum}}$ transcendunt, provenire unice ex iis terminis singulorum
coëfficientium $\gamma',\delta',\gamma'',\delta''$ etc., qui continent $\cos ny$ et
$\sin ny$ sive $\cos\mu nx$ et $\sin\mu nx$. Sint hi termini, in expressione pro
coëfficiente aliquo $\gamma^\lambda$ vel $\delta^\lambda$ $k\cos\mu nx+l\sin\mu nx$,
ex quibus producitur in $X''$
\[
\begin{aligned}
&(k\cos\mu nx+l\sin\mu nx)\cos\lambda x\quad\text{in casu primo}\\[5pt]
&\text{vel}\quad(k\cos\mu nx+l\sin\mu nx)\sin\lambda x\quad\text{in casu secundo}
\end{aligned}
\]
unde
\[
\begin{aligned}
&\tfrac12k\cos(\mu n+\lambda)x+\tfrac12l\sin(\mu n+\lambda)x\\[5pt]
&+\tfrac12k\cos(\mu n-\lambda)x+\tfrac12l\sin(\mu n-\lambda)x
\end{aligned}
\]
in casu primo atque
\[
\begin{aligned}
&-\tfrac12l\cos(\mu n+\lambda)x+\tfrac12k\sin(\mu n+\lambda)x\\[5pt]
&+\tfrac12l\cos(\mu n-\lambda)x-\tfrac12k\sin(\mu n-\lambda)x
\end{aligned}
\]
in casu secundo. Per art. 24 autem substituere oportet pro $\cos(\mu n+\lambda)x$
sive $\cos(\tfrac12\pi+\lambda)x$, $\cos\pi a\cos(\tfrac12\pi-\lambda)x+\sin\pi a\sin(\tfrac12\pi-\lambda)x$;
atque pro $\sin(\mu n+\lambda)x$ sive $\sin(\tfrac12\pi+\lambda)x$,
$\sin\pi a\cos(\tfrac12\pi-\lambda)x-\cos\pi a\sin(\tfrac12\pi-\lambda)x$.
Quare habebimus in casu primo
\[
\begin{aligned}
&\tfrac12(k+k\cos\pi a+l\sin\pi a)\cos(\tfrac12\pi-\lambda)x\\[5pt]
&+\tfrac12(l+k\sin\pi a-l\cos\pi a)\sin(\tfrac12\pi-\lambda)x
\end{aligned}
\]
in casu secundo autem
\[
\begin{aligned}
&\tfrac12(l-l\cos\pi a+k\sin\pi a)\cos(\tfrac12\pi-\lambda)x\\[5pt]
&-\tfrac12(k+l\sin\pi a+k\cos\pi a)\sin(\tfrac12\pi-\lambda)x
\end{aligned}
\]
\SourceLeaf{316}
Sed ex aequatione $k\sin\tfrac12\pi a=l\cos\tfrac12\pi a$, quae in utroque casu
inter $k$ et $l$ locum habere debet (art. 24), nullo negotio deducitur
\[
\begin{aligned}
&\tfrac12(k+k\cos\pi a+l\sin\pi a)=k\\[5pt]
&\tfrac12(l+k\sin\pi a-l\cos\pi a)=l
\end{aligned}
\]
Hinc tandem colligitur, pro iis partibus in $X''$, quae ex terminis ultimis
\[
\begin{aligned}
&k\cos ny+l\sin ny
\end{aligned}
\]
cuiusvis coëfficientis $\gamma^\lambda$ vel $\delta^\lambda$ producerentur,
designante $\lambda$ indicem maiorem quam 0, statim substitui posse
\[
\begin{aligned}
&k\cos(\tfrac12\pi-\lambda)x+l\sin(\tfrac12\pi-\lambda)x\\[5pt]
&\qquad\text{in casu priore vel}\\[5pt]
&l\cos(\tfrac12\pi-\lambda)x-k\sin(\tfrac12\pi-\lambda)x\\[5pt]
&\qquad\text{in casu posteriore}
\end{aligned}
\]
Hoc modo loco functionis $X''$, quae ad ordinem $\tfrac12\pi+m^{\mathrm{tum}}$
ascenderet, aliam $X''''$ nanciscimur, quae ordinem $\tfrac12\pi^{\mathrm{tum}}$
non egreditur. Nihil itaque superest, nisi ut coëfficientes cosinus et sinus arcus
$\tfrac12\pi x$, qui manifesto in hac functione iidem manserunt ut in $X''$, cum
coëfficientibus respondentibus in $X'''$ comparemus. Levi attentione autem
perspicitur, hosce terminos unice produci ex coëfficiente primo $\gamma$, et quidem
ex ultimis eius terminis; qui si per $k\cos ny+l\sin ny$ exhibentur, illi erunt
$k\cos\tfrac12\pi x+l\sin\tfrac12\pi x$, adeoque necessario identici cum
respondentibus in $X'''$. Hinc itaque colligitur, functionem $X''''$, quae per
praecepta modo tradita loco functionis $X''$ obtinetur, cum functione $X'''$
prorsus identicam fieri.

\InterpolationArticle{27}
Pro eo itaque casu, ubi multitudo valorum propositorum functionis $X$, periodum
integram formantium, numerus compositus est $=\pi=\mu\nu$, per partitionem illius
periodi in $\nu$ periodos $\mu$ terminorum eandem functionem cunctis valoribus
datis satisfacientem eruere in artt. 25, 26 didicimus, quae per applicationem
immediatam theoriae generalis ad periodum totam prodiret: illam vero methodum
calculi mechanici taedium magis minuere, praxis tentantem docebit. Nulla iam
amplius explicatione opus erit, quomodo illa partitio adhuc ulterius extendi et ad
eum casum applicari possit, ubi multitudo omnium valorum propositorum numerus
e tribus pluribusve factoribus compositus est, e. g. si numerus $\mu$
\SourceLeaf{317}
rursus esset, compositus, in quo casu manifesto quaevis periodus $\mu$ terminorum
in plures periodos minores subdividi potest. Ceterum ex pluribus aliis
annotationibus applicationi practicae inservientibus hic sequentes tantummodo
attingemus. Valores talium aggregatorum
\[
\begin{aligned}
&A\cos\lambda a+B\cos\lambda b+C\cos\lambda c\\[5pt]
&\qquad{}+D\cos\lambda d+\text{etc.}=p\\[5pt]
&A\sin\lambda a+B\sin\lambda b+C\sin\lambda c\\[5pt]
&\qquad{}+D\sin\lambda d+\text{etc.}=q
\end{aligned}
\]
(ubi ut in praecedentibus, supponimus $b-a=c-b=d\quad c\ \text{etc.}=\tfrac1\mu360^{\mathord{0}}=\Delta$)
saepius, siquidem nondum est $a=0$, commodius determinantur per formulas
\[
\begin{aligned}
&A+B\cos\lambda\Delta+C\cos2\lambda\Delta\\[5pt]
&\qquad{}+D\cos3\lambda\Delta+\text{etc.}=P\\[5pt]
&B\sin\lambda\Delta+C\sin2\lambda\Delta\\[5pt]
&\qquad{}+D\sin3\lambda\Delta+\text{etc.}=Q
\end{aligned}
\]
$p=P\cos a-Q\sin a$, $q=P\sin a+Q\cos a$, sive etiam faciendo
$\frac QP=\operatorname{tang}\varphi$, atque $\frac P{\cos\varphi}=\frac Q{\sin\varphi}=R$,
per has $p=R\cos(\varphi+a)$, $q=R\sin(\varphi+a)$.

\InterpolationArticle{28}
\textit{Exemplum.} In commercio literario a clar. barone \textsc{de Zach} edito,
vol. X p. 188 invenitur tabula, limitem borealem et australem zodiaci
\textit{Palladis} exhibens. Utriusque limites declinatio tamquam functio periodica
ascensionis rectae spectatur, quae in tabula illa per singulos quinque gradus
progreditur. Ad illustrationem disquisitionum praecedentium applicationem ad
limitem borealem hic faciemus, cuius itaque declinatio nobis erit $X$, ascensio
recta $x$. Excerpimus ex tabula illa periodum sequentem 12 terminorum
{\small\setlength{\arraycolsep}{3pt}\[
\begin{array}{r|r r l c r}
x&\multicolumn{5}{c}{X}\\ \hline
0^{\mathord{0}}&6^{\mathord{0}}&48'&\text{Bor.}&=+&408'\\
30&1&29&.\,.\,.\,.\,.\,.\,.&+&89\\
60&1&6&\text{Austr.}\,.\,.&-&66\\
90&0&10&\text{Bor.}\,.\,.\,.&+&10\\
120&5&38&.\,.\,.\,.\,.\,.\,.&+&338\\
150&13&27&.\,.\,.\,.\,.\,.\,.&+&807\\
180&20&38&.\,.\,.\,.\,.\,.\,.&+&1238\\
210&25&11&.\,.\,.\,.\,.\,.\,.&+&1511\\
240&26&23&.\,.\,.\,.\,.\,.\,.&+&1583\\
270&24&22&.\,.\,.\,.\,.\,.\,.&+&1462\\
300&19&43&.\,.\,.\,.\,.\,.\,.&+&1183\\
330&13&24&.\,.\,.\,.\,.\,.\,.&+&804
\end{array}
\]}
\SourceLeaf{318}
Distribuamus hanc periodum primo in tres periodos quaternorum terminorum
{\small\setlength{\arraycolsep}{3pt}\[
\begin{array}{r r|r r|r r}
a=0^{\mathord{0}}&A=+408&a'=30^{\mathord{0}}&A'=+89&a''=60^{\mathord{0}}&A''=-66\\
b=90^{\mathord{0}}&B=+10&b'=120^{\mathord{0}}&B'=+338&b''=150^{\mathord{0}}&B''=+807\\
c=180^{\mathord{0}}&C=+1238&c'=210^{\mathord{0}}&C'=+1511&c''=240^{\mathord{0}}&C''=+1583\\
d=270^{\mathord{0}}&D=+1462&d'=300^{\mathord{0}}&D'=+1183&d''=330^{\mathord{0}}&D''=+804
\end{array}
\]}
In formula
\[
\begin{aligned}
&X'={}\gamma+\gamma'\cos x+\gamma''\cos2x\\[5pt]
&+\delta'\sin x+\delta''\sin2x
\end{aligned}
\]
fit itaque
\[
\begin{aligned}
&\gamma=\tfrac14(A+B+C+D)=779'{,}5\\[5pt]
&\gamma'=\tfrac12(A\cos a+B\cos b+C\cos c+D\cos d)\\[5pt]
&=\tfrac12(A-C)=-415'{,}0\\[5pt]
&\delta'=\tfrac12(A\sin a+B\sin b+C\sin c+D\sin d)\\[5pt]
&=\tfrac12(B-D)=-726'{,}0\\[5pt]
&\gamma''=\tfrac14(A\cos2a+B\cos2b+C\cos2c+D\cos2d)\\[5pt]
&=\tfrac14(A-B+C-D)\\[5pt]
&=+43'{,}5\\[5pt]
&\delta''=\tfrac14(A\sin2a+B\sin2b+C\sin2c+D\sin2d)=0
\end{aligned}
\]
et similiter pro periodo secunda ac tertia. Hoc modo emergit
{\small\setlength{\arraycolsep}{3pt}\[
\begin{array}{l|r|r|r|r|r|r}
\text{Pro periodo}&\text{ubi }y=4x&\gamma&\gamma'&\delta'&\gamma''&\delta''\\ \hline
\text{prima}&0^{\mathord{0}}&+779{,}5&-415{,}0&-726{,}0&+43{,}5&0\\
\text{secunda}&120^{\mathord{0}}&+780{,}2&-404{,}5&-721{,}4&+9{,}9&+17{,}1\\
\text{tertia}&240^{\mathord{0}}&+782{,}0&-413{,}5&-713{,}3&+11{,}7&-20{,}3
\end{array}
\]}
Hic porro exhibetur $\gamma$ per formulam
\[
\begin{aligned}
&\tfrac13(779'{,}5+780{,}2+782{,}0)\\[5pt]
&+\tfrac23(779{,}5+780{,}2\cos120^{\mathord{0}}+782{,}0\cos240^{\mathord{0}})\cos4x\\[5pt]
&+\tfrac23(780{,}2\sin120^{\mathord{0}}+782{,}0\sin240^{\mathord{0}})\sin4x
\end{aligned}
\]
sive per
\[
\begin{aligned}
&780{,}6-1{,}1\cos4x-1{,}0\sin4x,\quad\text{et perinde}
\end{aligned}
\]
\[
\begin{aligned}
&\gamma'\quad\text{per}\quad-411{,}0-4{,}0\cos4x+5{,}2\sin4x\\[5pt]
&\delta'\quad\text{per}\quad-720{,}2-5{,}8\cos4x-4{,}7\sin4x\\[5pt]
&\gamma''\quad\text{per}\quad+21{,}7+21{,}8\cos4x-1{,}1\sin4x\\[5pt]
&\delta''\quad\text{per}\quad-1{,}1+1{,}1\cos4x+21{,}6\sin4x
\end{aligned}
\]
Quibus valoribus in $X'$ substitutis prodit formula 12 valores propositos exhibens
\SourceLeaf{319}
\[
\begin{aligned}
&+780{,}6\\[5pt]
&-411{,}0\cos x-720{,}2\sin x\\[5pt]
&+43{,}4\cos2x-2{,}2\sin2x\\[5pt]
&-4{,}3\cos3x+5{,}5\sin3x\\[5pt]
&-1{,}1\cos4x-1.0\sin4x\\[5pt]
&+0{,}3\cos5x-0{,}3\sin5x\\[5pt]
&+0{,}1\cos6x
\end{aligned}
\]
Distribuamus secundo eandem periodum in quatuor periodos ternorum terminorum, quibus singulis itaque per formulam primi ordinis
\[
\begin{aligned}
&X'=\gamma+\gamma'\cos x+\delta'\sin x
\end{aligned}
\]
satisfaciendum erit. Hic invenitur
\begin{center}
\begin{tabular}{r|c|r|r|r}
Pro periodo & ubi $y=3x$ & $\gamma$ & $\gamma'$ & $\delta'$\\\hline
prima & $0^{0}$ & $+776{,}3$ & $-368{,}3$ & $-718{,}8$\\
secunda & $90^{0}$ & $+786{,}0$ & $-414{,}5$ & $-676{,}0$\\
tertia & $180^{0}$ & $+785{,}0$ & $-453{,}0$ & $-721{,}1$\\
quarta & $270^{0}$ & $+775{,}0$ & $-408{,}2$ & $-765{,}0$
\end{tabular}
\end{center}
unde deducuntur formulae, sub quibus $\gamma,\gamma',\delta'$ exhibentur sequentes:
\[
\begin{aligned}
&\gamma\quad.\,.\,.\,.\quad+780{,}6-4{,}3\cos3x+5{,}5\sin3x\\[5pt]
&\qquad{}+0{,}1\cos6x\\[5pt]
&\gamma'\quad.\,.\,.\,.\quad-411{,}0+42{,}3\cos3x-3{,}2\sin3x\\[5pt]
&\qquad{}+0{,}3\cos6x\\[5pt]
&\delta'\quad.\,.\,.\,.\quad-720{,}2+1{,}2\cos3x+44{,}5\sin3x\\[5pt]
&\qquad{}+0{,}3\cos6x
\end{aligned}
\]
His valoribus in $\gamma+\gamma'\cos x+\delta'\sin x$ substitutis, partibusque ultimis in $\gamma,\gamma',\delta'$, quae $\cos6x$ continent, secundum praecepta art. 26 tractatis, prorsus eadem formula eruitur, ad quam supra pervenimus.
\InterpolationArticle{29}
Supersunt casus, ubi in functione $X$ vel cosinus vel sinus soli adsunt, quos in artt. 16, 17 generaliter tractavimus. Supponemus, ut in praecc., datam esse periodum $\mu$ terminorum, puta $X=A,B,C,D$ etc. pro $x=a,b,c,d$ etc. Duo vero casus hic probe distinguendi sunt. Aliter enim tractari debet

$1^{\mathrm{mo}}$ casus, ubi $a$ est complementum alicuius valorum sequentium $b,c,d$ etc. ad $360^{0}$ vel ad $0$ vel ad multiplum peripheriae. Sit hic terminus $a+\frac{k}{\mu}360^{0}$,
\SourceLeaf{320}
atque $a+\frac{k}{\mu}360^{0}+a=l\times360^{0}$, sive $a=(\mu l-k)\frac{180}{\mu}$. Quare in hoc casu $a$ per $\frac{180}{\mu}$ divisibilis erit, adeoque $\sin\mu a=0$; porro habetur $b=(\mu l-k+2)\frac{180^{0}}{\mu}$, atque $\sin\mu b=0$, et perinde de sequentibus $c,d$ etc. Supponendo itaque, $a$ esse minorem quam $\frac{360^{0}}{\mu}$, sed non negativum (quod permissum est, quum ab omnibus $a,b,c,d$ etc. multiplum proxime minus totius peripheriae subtrahere et ex residuis minimum, quod proprietate illa praeditum erit pro periodi initio accipere liceat), erit vel $a=0$, atque $b=\frac1\mu360^{0}$, $c=\frac2\mu360^{0}$, $d=\frac3\mu360^{0}$ etc.; vel $a=\frac1\mu180^{0}$, $b=\frac3\mu180^{0}$, $c=\frac5\mu180^{0}$, $d=\frac7\mu180^{0}$ etc.

aliter $2^{\mathrm{do}}$ casus, ubi non est complementum ullius valorum sequentium $b,c,d$ etc. ad $360^{0}$ vel $0$ vel multiplum peripheriae. Hic $a$ non erit divisibilis per $\frac{180^{0}}{\mu}$, adeoque $\sin\mu a=\sin\mu b=\sin\mu c=\sin\mu d$ etc. non $=0$.

In casu posteriore methodus generalis art. 16 statim applicari potest. In priore autem ex valoribus $a,b,c,d$ etc. (quos infra $360^{0}$ reductos esse supponimus), omnes antea reiicere oportet, qui sunt maiores quam $180^{0}$, quippe quorum complementa ad $360^{0}$ inter reliquos reperiuntur; insuperque, quando adsunt, valores $0$ et $180^{0}$, si in functione $X$ sinus soli occurrunt: sed quoniam applicatio methodi generalis hac ratione minus concinna evaderet, hunc casum alio modo infra absolvemus. Initium iam ab illo casu faciemus, ubi $\sin\mu a$ non est $=0$.
\InterpolationArticle{30}
Sit primo $X$ functio formae
\[
\begin{aligned}
&\alpha+\alpha'\cos x+\alpha''\cos2x+\alpha'''\cos3x+.\,.\,.+\alpha^{n}\cos nx
\end{aligned}
\]
unde esse debebit $\mu=n+1$. Hic habetur per lemma secundum art. 19
\[
\begin{aligned}
&(\cos t-\cos b)(\cos t-\cos c)(\cos t-\cos d)\ \text{etc.}
\\[5pt]
&=\frac1{2^{\mu-1}}\times\frac{\cos\mu t-\cos\mu a}{\cos t-\cos a}\\[5pt]
&=\frac1{2^{\mu-1}\sin a}\{\sin\mu a+2\sin(\mu-1)a\cos t\\[5pt]
&\qquad{}+2\sin(\mu-2)a\cos2t+\text{etc.}+2\sin a\cos(\mu-1)t\}
\end{aligned}
\]
Facile scilicet confirmatur, productum ex aggregato secundae partis huius aequationis per $\cos t-\cos a$ fieri $=\sin a(\cos\mu t-\cos\mu a)$. Quare fit quoque, scribendo $a$ pro $t$
\[
\begin{aligned}
&(\cos a-\cos b)(\cos a-\cos c)(\cos a-\cos d)\ \text{etc.}
\\[5pt]
&=\frac1{2^{\mu-1}\sin a}\{\sin\mu a+2\sin(\mu-1)a\cos a\\[5pt]
&\qquad+2\sin(\mu-2)a\cos2a+\text{etc.}+2\sin a\cos(\mu-1)a\}
\end{aligned}
\]
Aggregatum in hac expressione fit $=\mu\sin\mu a$, ut inde manifestum est, quod
\SourceLeaf{321}
\[
\begin{aligned}
&\text{terminus secundus}=\sin\mu a+\sin(\mu-2)a\\[5pt]
&\text{ultimus}=\sin\mu a-\sin(\mu-2)a\\[5pt]
&\text{porro terminus tertius}=\sin\mu a+\sin(\mu-4)a\\[5pt]
&\text{terminus penultimus}=\sin\mu a-\sin(\mu-4)a\\[5pt]
&\text{etc.}
\end{aligned}
\]
Hinc colligitur, coëfficientem ipsius $A$ in formula art. 16 pro $T$ fieri
\[
\begin{aligned}
&\frac1{\mu\sin\mu a}\{\sin\mu a+2\sin(\mu-1)a\cos t\\[5pt]
&\qquad{}+2\sin(\mu-2)a\cos2t+\text{etc.}+2\sin a\sin(\mu-1)t\}
\end{aligned}
\]
Prorsus similes expressiones proveniunt pro coëfficientibus ipsorum $B,C,D$ etc., mutando tantummodo $a$ in $b,c,d$ etc. Quamobrem statuendo
\[
\begin{aligned}
&\varepsilon=\frac1\mu(A+B+C+D+\text{etc.})\\[5pt]
&\varepsilon'=\frac2{\mu\sin\mu a}\{A\sin(\mu-1)a+B\sin(\mu-1)b\\[5pt]
&\qquad{}+C\sin(\mu-1)c+D\sin(\mu-1)d+\text{etc.}\}\\[5pt]
&\varepsilon''=\frac2{\mu\sin\mu a}\{A\sin(\mu-2)a+B\sin(\mu-2)b\\[5pt]
&\qquad{}+C\sin(\mu-2)c+D\sin(\mu-2)d+\text{etc.}\}\\[5pt]
&\varepsilon'''=\frac2{\mu\sin\mu a}\{A\sin(\mu-3)a+B\sin(\mu-3)b\\[5pt]
&\qquad{}+C\sin(\mu-3)c+D\sin(\mu-3)d+\text{etc.}\}\\[5pt]
&\text{etc.}\\[5pt]
&\varepsilon^{n}=\frac2{\mu\sin\mu a}(A\sin a+B\sin b\\[5pt]
&\qquad{}+C\sin c+D\sin d+\text{etc.})
\end{aligned}
\]
erit
\[
\begin{aligned}
&T=\varepsilon+\varepsilon'\cos t+\varepsilon''\cos2t+\varepsilon'''\cos3t+\text{etc.}\\[5pt]
&\qquad{}+\varepsilon^{n}\cos nt
\end{aligned}
\]
Quum haec formula generaliter pro valore quocunque ipsius $t$ valere debeat, necessario cum $X$ identica fiet, mutata $t$ in $x$, adeoque coëfficientes $\varepsilon,\varepsilon',\varepsilon''$ etc. ipsis $\alpha',\alpha',\alpha''$ etc. resp. aequales.
\InterpolationArticle{31}
Quum per praecepta supra explicata functio $X'$ formae
\[
\begin{aligned}
&\gamma+\gamma'\cos x+\gamma''\cos2x+.\,.\,.+\gamma^{m}\cos mx\\[5pt]
&\quad+\delta'\sin x+\delta''\sin2x+.\,.\,.\,.+\delta^{m}\sin mx
\end{aligned}
\]
quae omnibus $\mu$ valoribus propositis satisfaciat, et in qua sit $m=\frac12\mu-\frac12$ vel $\frac12\mu$, prout $\mu$ impar est vel par: operae pretium est, hanc functionem cum functione modo inventa
\[
\begin{aligned}
&\varepsilon+\varepsilon'\cos x+\varepsilon''\cos2x+\varepsilon'''\cos3x+.\,.\,.\\[5pt]
&\qquad{}+\varepsilon^{\mu-1}\cos(\mu-1)x
\end{aligned}
\]
\endgroup
